% !TeX root = ../大学生数学竞赛概率论与数理统计试题选编与专题训练.tex
% 由主文档通过 \input 载入；请编译主文档。
\chapter*{编写说明与阅读路线}
\addcontentsline{toc}{chapter}{编写说明与阅读路线}
本书面向已修读微积分和基础线性代数、准备进一步学习概率论与数理统计竞赛问题的本科生。全国大学生数学竞赛（CMC）的概率相关内容主要出现在数学类高年级决赛；它与有独立“概率与统计”科目的丘成桐大学生数学竞赛属于不同赛事。本书分别编排这两类题源，保留届次、年份、组别、题号和可核查入口，不把丘赛试题称为 CMC 试题。

全书由竞赛真题选讲、专题方法与课程衔接训练组成。竞赛部分收录十三道 CMC 概率相关题与72道丘赛选题，共85个来源题目条目；其中丘赛2019年原题4和2025年原题1只选录明确标注的部分子问。训练部分整理本项目既有《概率论与数理统计习题集》中的264道课程题及其解答，涵盖条件概率、分布、矩、极限定理、参数估计、检验和回归。这264题是课程训练，并非竞赛真题。题面和解答分列时仍按同一道题计数，题目子问也不重复计为独立真题。全书页数包含解答、方法、索引和课程训练，不能解读成同样页数的纯竞赛原卷。

CMC 题面和基础题解来自 Jacobi Lie 的\href{https://github.com/lzxcr/cmc-mathematics/tree/64916fa1fdbc3e444a652eca204d5669058ee321}{《全国大学生数学竞赛题解》固定版本（64916fa）}，来源层级为公开题解转录。原作者的题解、评注、方法说明采用\href{https://creativecommons.org/licenses/by-sa/4.0/}{CC BY-SA 4.0}，本书相应改编部分保留署名并沿用该许可，新增的排版和条件说明不代表原作者审定。试题本身及其他第三方材料仍归相应权利人。

丘赛题源为\href{https://yau-contest.com/lists-jxxg.html}{组委会官方“相关下载”}中“2010--2026历年笔试真题”压缩包，\href{https://yau-contest.com/index/reldown?id=9}{官方原卷下载}。本书实际选题覆盖2012--2026年，将数学条件译写为中文，解答由编订者推导；选取某题的部分子问时会明确标记。官方题包中的2020--2022年附有参考解答，2022年原题5的边界情形另作了核查说明。本书未将自己的证明称为官方答案。题面记号经过统一重写；有争议的题意、遗漏条件及不成立的结论在题后说明。

课程训练沿用已有中文开放题解及 MIT OCW 18.05 材料的来源标记，完整署名、许可、原题定位和校订记录见训练部分书末。不同来源分部独立保留原有许可；本书的集合编排不以单一许可覆盖或改变所有材料。课程题在本书中重新编号，题目、答案和来源索引通过交叉引用对应。

建议先阅读与已学课程对应的竞赛题，在解答中遇到条件期望、紧性、截断、充分性或随机化推断等概念时，转入方法部分；基础计算不熟练时再使用训练部分。布朗运动、测度变换和一般马尔可夫链属于进阶内容，不能将其误解为普通工科概率统计课程的统一要求。本书采用 A4 开本，排版沿用项目共享样式。

\section*{记号与参数约定}
$\PP$、$\E$、$\Var$、$\Cov$ 分别表示概率、期望、方差和协方差。$\operatorname{Exp}(\lambda)$ 除非题目明确采用均值参数，否则使用率参数，期望为 $1/\lambda$。$\operatorname{Gamma}(a,\lambda)$ 使用形状 $a$ 和率 $\lambda$。$\Phi$ 与 $\phi$ 是标准正态分布函数和密度。$X_{(1)}\le\cdots\le X_{(n)}$ 是次序统计量。$\mathbf1_A$ 表示事件 $A$ 的示性函数；$\xrightarrow{P}$、$\xrightarrow{d}$、$\xrightarrow{\mathrm{a.s.}}$ 分别表示依概率、依分布和几乎处处收敛。

样本变量用大写，观测值用小写；未知参数通常视为固定量。贝叶斯模型会另行指定先验。$z_q,t_{\nu;q},\chi^2_{\nu;q},F_{a,b;q}$ 均表示下侧累计概率为 $q$ 的分位数。统计推断中的“精确”指在题设模型内成立；渐近公式不承诺有限样本精确覆盖。所有积分换序、极限交换和无偏性结论均须检查所列可积性或抽样条件。

MLE 表示\emph{极大似然估计}（maximum likelihood estimator），UMVUE 表示\emph{一致最小方差无偏估计}（uniformly minimum variance unbiased estimator）。后者中的“一致”指对参数空间中每个参数值都具有最小方差，不等同于随样本量增长的一致性（consistency）。本书讨论 UMVUE 时采用有限方差估计量的通常约定。$O_P(1)$ 表示随机变量列依概率有界，$o_P(1)$ 表示依概率趋零；更一般的随机阶符号通过相应的确定性尺度归一化定义。
\tableofcontents
\clearpage






